The analysis of complex systems has long sought intermediate constructs that are simultaneously mathematically tractable and ontologically meaningful. Recurrence-based methods provide powerful descriptors of local persistence and determinism, yet translating these descriptors into coherent units of analysis that respect both the mathematics of nonlinear dynamics and the structure of temporal becoming remains an open challenge.

Within the RECD + Ontología del Presente program, the three-layer framework distinguishes (i) local intensification of persistence (Layer 1), (ii) sustained coherence of persistence scales across modules (Layer 2), and (iii) global structural reorganization (Layer 3). The central proposal of this paper is that the \textbf{Joint Episode} — the co-occurrence of Layer 1 within a sustained Layer 2 window — constitutes a primitive ontological-temporal unit.

This paper develops the Joint Episode rigorously on three fronts:
\begin{itemize}
    \item Formal mathematical definition and associated metrics (Joint Score, Confidence Score, duration, intensity).
    \item Derivation of structural properties and relations to anti-synchronization, memory, and layer transitions.
    \item Validation in controlled synthetic systems (coupled logistic maps, Lorenz attractor) together with an explicit ontological interpretation.
\end{itemize}

The emphasis is technical and mathematical, with the ontological dimension articulated primarily in its dedicated section while informing the formal development throughout.
